\documentclass[10pt]{article}
%% Shared preamble for the 9.520 Oct 2 literature reviews.
\usepackage[T1]{fontenc}
\usepackage[utf8]{inputenc}
\usepackage{lmodern}
\usepackage[margin=0.75in]{geometry}
\usepackage{amsmath,amssymb}
\usepackage{microtype}
\usepackage{enumitem}
\usepackage[dvipsnames]{xcolor}
\usepackage[backend=bibtex,style=authoryear,natbib=true,maxcitenames=2,maxbibnames=4,minbibnames=3,uniquelist=false,uniquename=false,giveninits=true,doi=false,isbn=false,url=false,eprint=true]{biblatex}
\addbibresource{refs.bib}
\usepackage[colorlinks=true,linkcolor=MidnightBlue,citecolor=MidnightBlue,urlcolor=MidnightBlue]{hyperref}
\usepackage[capitalise,nameinlink]{cleveref}
\setlength{\parskip}{2pt}
\setlength{\parindent}{0pt}
\setlist{nosep,leftmargin=*}

%% Tags for the annotated entries: [F] foundational, [R] 2024-26, [P] preprint or workshop.
\newcommand{\tagF}{\textsf{\textcolor{MidnightBlue}{[F]}}}
\newcommand{\tagR}{\textsf{\textcolor{ForestGreen}{[R]}}}
\newcommand{\tagP}{\textsf{\textcolor{BrickRed}{[P]}}}

%% One annotated bibliography entry: \entry{key}{tags}{annotation}
\DeclareBibliographyCategory{annotated}
\newcommand{\entry}[3]{%
  \addtocategory{annotated}{#1}%
  \par\noindent\hangindent=1.2em\hangafter=1
  {\small\fullcite{#1}}~#2\par
  {\leftskip=1.2em #3\par}\vspace{2pt}}

%% Group synthesis paragraph opening each explanation.
\newcommand{\synth}[1]{\par\noindent\emph{#1}\par\vspace{2pt}}

%% Works cited in the text but not annotated, printed compactly at the end.
\newcommand{\otherrefs}{%
  {\footnotesize\setlength{\bibitemsep}{0pt}%
  \printbibliography[notcategory=annotated,heading=subbibliography,title={Other references}]}}


\title{\vspace{-1.6em}Optimizer-Controlled Implicit Bias and Functional Complexity\\[2pt]
\large Literature review and implications (9.520, Fall 2026)}
\author{Juntang Wang \and V\'ictor Conchello}
\date{October 2, 2026}

%% Split: Juntang wrote Sec. 1, 2.1-2.3, 2.5, the key-papers draft and the implications draft.
%% Victor writes Sec. 2.4 (Adam and Muon) and edits Sec. 3-4.
%% Page budget: the document must stay at 4 pages, so Sec. 2.4 gets at most ~0.5 page
%% (4-5 entries of ~45 words).

\begin{document}
\maketitle
\vspace{-1.5em}

%% =================================================================
%% Question
%% =================================================================
\section{The question and why it is hard}
\label{sec:question}

Overparameterized networks have many parameter settings with the same loss; the optimizer picks
which one training returns. The project asks whether SGD, Adam and Muon, \emph{at matched loss},
select functions of different complexity, curvature or robustness, and whether swapping optimizers
mid-training moves the solution toward the new optimizer's preferred one. Three obstacles: theory
states implicit bias mostly for linear, diagonal or infinitely wide models, as parameter norms that
differ by optimizer; experiments report test accuracy, which cannot separate risk-equivalent
predictors; and curvature depends on parameterization, so ``flatter'' need not mean ``simpler''.
Tags: \tagF{} foundational, \tagR{} 2024--26 top venue, \tagP{} preprint or workshop.

%% =================================================================
%% Literature
%% =================================================================
\section{Annotated literature}
\label{sec:lit}

\subsection{Foundations: the optimizer's geometry selects the solution}
\label{sec:lit-foundations}

\synth{Implicit bias belongs to the optimizer and the parameterization together, not to the loss. In
linear models a first-order method acts as steepest descent in some norm and converges to the
max-margin solution in that norm, but slowly: GD reaches the $\ell_2$ margin on separable data only
at a logarithmic rate \citep{soudry2018separable}, so a finite-time comparison may be far from the
limit.}

\entry{gunasekar2018geometry}{\tagF}{For mirror descent, natural gradient and steepest descent on
underdetermined regression and separable linear classification, asks when the limit depends only on
the geometry, not on step size or momentum. For steepest descent on separable data it is the
max-margin solution in the chosen norm. Later results for Adam ($\ell_\infty$) and Muon (spectral
norm) follow this template.}

\entry{gunasekar2018conv}{\tagF}{GD on full-width linear convolutional networks of depth $L$ converges
to a predictor tied to an $\ell_{2/L}$ bridge penalty in the frequency domain; fully connected linear
networks reach the hard-margin SVM at any depth. Same optimizer, different function, so any
optimizer comparison is also an architecture comparison.}

\entry{chizat2020logistic}{\tagF}{For infinitely wide two-layer homogeneous networks, gradient flow on
the logistic loss converges to a max-margin classifier in a non-Hilbertian function space that
adapts to hidden low-dimensional structure; training only the last layer gives a kernel SVM without
that adaptivity. The bias is stated for \emph{functions}, the level at which we compare optimizers.}

\entry{wilson2017marginal}{\tagF}{On constructed separable data, GD and SGD reach zero test error
while AdaGrad, RMSProp and Adam approach chance; on deep-learning benchmarks adaptive methods
generalize worse even when they train better. The claim we test, that Adam finds different
solutions, starts here.}

\subsection{Measuring functional complexity}
\label{sec:lit-metrics}

\synth{We need two independent function metrics. These are defined on the input-output map, so unlike
sharpness or weight norms they do not depend on parameterization. Region counts say how many pieces
the function has; Jacobian norms say how steep they are. Both change at constant training loss:
regions drift from the training points toward the decision boundary long after interpolation
\citep{humayun2024grok}, so a matched-loss comparison must also fix training time.}

\entry{hanin2019regions}{\tagF}{At initialization, the number of linear regions of a ReLU network
along a line grows linearly in the number of neurons, far below the exponential bound, and stays far
below it after training. Region counts are computable and typically small.}

\entry{novak2018sensitivity}{\tagF}{Across thousands of trained models, the input-output Jacobian norm
is small near the training data and correlates with generalization across architectures, optimizers
and datasets. It measures slope, complementing region counts.}

\entry{li2025regions}{\tagR}{Counts connected input regions with the same predicted label, a
parameterization-free measure of implicit bias. Small counts go with simple boundaries and good
generalization; larger learning rates and smaller batches give smaller counts, with theory for the
learning-rate effect. It compares hyperparameters, not optimizers.}

\subsection{The GD and SGD side: step size and noise select simpler functions}
\label{sec:lit-sgd}

\synth{The best-developed side, and the closest to the course's training-dynamics line. Step size
acts through stability, mini-batch noise through an implicit penalty, scaled by $\eta$ or $\eta/B$.
What they buy depends on the model: smoothness in low dimension, a sparse input support, low rank.}

\entry{qiao2024stable}{\tagR}{\citet{mulayoff2021stability} showed that minima GD can stably reach at
step $\eta$ in univariate two-layer ReLU networks have weighted total variation of $f'$ at most about
$1/\eta$. This paper drops interpolation: with label noise $\sigma$, stable minima have weighted total
variation at most $1/\eta-1/2+\widetilde O(\sigma+\sqrt{\mathrm{MSE}})$, giving a near-optimal
$\widetilde O(n^{-4/5})$ rate. Univariate inputs only.}

\entry{liang2025shattering}{\tagR}{With multivariate inputs, flatness still implies generalization,
but the rates worsen exponentially with dimension: flat solutions can ``shatter'' the data with
rarely active, large neurons. Low-norm solutions avoid this, so stability alone cannot explain
generalization in high dimension.}

\entry{smith2021origin}{\tagF}{For small finite $\eta$, the mean SGD iterate with random shuffling
follows gradient flow on a modified loss that penalizes mini-batch gradient norms, with strength
proportional to $\eta/B$ for small batches. ``SGD equals GD plus a step on a regularizer'' becomes a
theorem, with the sign of the batch-size effect.}

\entry{beneventano2024support}{\tagP}{From the instructors' group. Mini-batch SGD learns the target's
support in the first layer, driving weights on irrelevant inputs to zero through a second-order
implicit regularizer proportional to $\eta/B$; full-batch GD needs an explicit penalty to do the
same. Relatedly, SGD with weight decay favours low-rank weights \citep{galanti2022rank} and can jump
to lower-rank minima but never back \citep{wang2024onewayjumps}.}

\subsection{The adaptive side: Adam and Muon}
\label{sec:lit-adaptive}

%% VICTOR: at most ~0.5 page. Suggested entries: Zhang-Zou-Cao 2024 (NeurIPS), Xie-Li 2024 (ICML),
%% Fan-Schmidt-Thrampoulidis 2025 (NeurIPS), Beneventano-Abdelmoneum-Poggio 2026 (preprint), plus
%% Liu et al. 2025 for Muon at scale. Once added, cite Xie-Li and Fan et al. in implication 1.
\synth{[V\'ictor's section: Adam and sign descent converge toward $\ell_\infty$ max-margin solutions,
Muon and spectral descent toward the spectral-norm max margin, and their implicit regularizers can
penalize or anti-penalize gradient norms.]}

\subsection{Matched-loss comparisons and optimizer swaps}
\label{sec:lit-matched}

\synth{These show the function can change while the error does not. None compares SGD, Adam and Muon
at matched loss with two independent function metrics, and the one swap study reports only accuracy.
That is the gap.}

\entry{cortesi2026same}{\tagP}{From the course staff. In S\&P~500 volatility forecasting, models with
identical test loss learn different functions; the optimizer reshapes response profiles and temporal
dependence, with about $3\times$ spread in portfolio turnover at similar Sharpe ratio. We extend this
matched-loss design beyond one noisy domain.}

\entry{vasudeva2025rich}{\tagR}{For two-layer ReLU networks on Gaussian data, proves GD has a
simplicity bias toward a linear boundary with suboptimal margin, while Adam learns richer features and
a nonlinear boundary closer to Bayes-optimal; empirically Adam then generalizes better under spurious
correlations. Here richness helps, the opposite of the usual lesson from
\citet{wilson2017marginal}.}

\entry{keskar2017swats}{\tagF\,\tagP}{SWATS starts with Adam and switches to SGD when a
projection-based trigger fires, closing the generalization gap on most CIFAR, Tiny-ImageNet and
language-modeling tasks tried. It is the existing swap intervention, but it reports only accuracy,
so it cannot say whether the swap changes the learned function.}

%% =================================================================
%% Key papers
%% =================================================================
\section{Three key papers}
\label{sec:key}

\textbf{\citet{wilson2017marginal}} (NeurIPS; about 1{,}200 citations on Semantic Scholar) started the
question of whether optimizers differ in what they learn, not just how fast. Its construction proves
adaptive and non-adaptive methods can reach opposite test errors on the same data, and its
experiments made ``SGD generalizes better'' a belief later work has had to qualify. Our matched-loss
comparisons re-test that claim with modern optimizers.

\textbf{\citet{gunasekar2018conv}} (NeurIPS; about 490 citations; listed in the project description)
shows implicit bias depends on parameterization as much as on the optimizer, and characterizes it as
a function-space penalty in a nontrivial model, part of the Srebro and Soudry work that made implicit
bias a theorem. Design rule for us: hold the architecture fixed, and do not assume results transfer
across architectures.

\textbf{\citet{vasudeva2025rich}} (NeurIPS 2025) is the most recent top-venue theorem comparing Adam
and GD in function space for nonlinear networks. It names a mechanism (GD's simplicity bias versus
Adam's richer features), predicts boundary shape, and shows ``simpler'' is not always better. It is
the first hypothesis to test on our matched-loss model bank.

%% =================================================================
%% Implications
%% =================================================================
\section{Implications for theory and practice}
\label{sec:implications}

\begin{enumerate}[leftmargin=1.3em,itemsep=4pt]
\item \textbf{Comparing complexity across optimizers is circular unless the metric is neutral.} Each
optimizer favours small values of the norm it descends in \citep{gunasekar2018geometry}: $\ell_2$
for GD, $\ell_\infty$ for Adam-like methods, spectral for Muon (\cref{sec:lit-adaptive}). Measuring
complexity in one of these norms favours the matching optimizer by construction. Use metrics defined
on the function, such as region counts \citep{hanin2019regions,li2025regions} and input-Jacobian norms
\citep{novak2018sensitivity}, and require the two to agree.

\item \textbf{Theory: stability explains the GD and SGD side only in low dimension, so structure must
do the rest.} Minima stability gives sharp function-space bounds in one dimension
\citep{mulayoff2021stability,qiao2024stable}, but the rates degrade exponentially with input
dimension \citep{liang2025shattering}. Explanations that survive must use structure: input support
\citep{beneventano2024support}, rank \citep{galanti2022rank,wang2024onewayjumps}, or more generally
compositional sparsity, where the relevant dimension is that of each constituent, not the input.
Prediction: optimizer differences in complexity should be largest on compositionally sparse targets,
where the $\ell_2$, $\ell_\infty$ and spectral geometries disagree most about which constituents to
use.

\item \textbf{Practice: matched test loss is not enough to choose an optimizer.} With risk-equivalent
predictors, the optimizer still changes robustness, response profiles and downstream decisions
\citep{cortesi2026same,vasudeva2025rich}. Report at least one function-level metric next to the loss,
especially in noisy domains. A swap at matched loss, as in SWATS \citep{keskar2017swats} but scored
with function metrics, is a cheap test: if the function moves, the optimizer is a modeling choice,
not a tuning detail.

\item \textbf{Open question: is the bias set early or applied throughout training?} The break-even
picture \citep{jastrzebski2020breakeven} and the EoS literature suggest hyperparameters fix a
curvature regime early. Swapping at different times, scored with two neutral metrics, separates a
bias that acts by choosing the basin early from one that keeps acting inside it. This is the causal
form of the project's question, and it connects to the instructors' work on stability with momentum
\citep{andreyev2026momentum} and under Muon (\cref{sec:lit-adaptive}).
\end{enumerate}

\otherrefs

\vspace{2pt}
{\footnotesize\textbf{Use of generative AI.} We used an AI assistant for literature discovery,
checking citation metadata against arXiv and proceedings pages, and drafting prose. We read the
primary sources for every claim above, verified all citations, and take responsibility for the
content.}

\end{document}
